\section{The adapted transfer protocol of the original submission}
\label{sec:supp_adapted}

The cross dataset numbers of the original submission were not produced by the no adaptation
protocol the text described. The script that produced them resumed the corrector trained on PKU37
and trained it further on the target dataset in a leave one subject out design. We rerun that
protocol here with the selected NAFNet checkpoint so that the gap between the two protocols can be
read directly, and we describe it exactly.

\subsection{Protocol}
Each fold holds out one subject, seventeen of Duke2013 or fifteen of Duke17 subjects form the
training set, and the held out subject is evaluated once. Training resumes from the PKU37 checkpoint
of the tissue bounded gate setting at seed 0 and runs for fifteen epochs with batch size two. The
backbone weights and its normalisation statistics are frozen. The wrapper and the cooperation head
that reads the backbone features are trainable. The learning rates are $2{\times}10^{-5}$,
$1{\times}10^{-4}$ and $6{\times}10^{-5}$ for the corrector, potential and negotiator, five to ten
times below the PKU37 schedule of Table~\ref{tab:constants}, the fidelity term enters at epoch five,
and an L2 penalty of weight 0.1 holds every trainable weight near its PKU37 value, which is elastic
weight consolidation without the Fisher weighting. A background correction variance penalty of
weight 10 is also active. The checkpoint kept for each fold is the epoch with the highest validation
score, and the validation image of a fold is its held out subject. The adapted numbers are therefore
selected on the subject they are evaluated on and should be read as an upper bound for this
protocol. The same was true of the original submission.

\subsection{Result}
Table~\ref{tab:supp_adapted} gives both target sets, next to the no adaptation result of the same
checkpoint from the main paper.

\begin{table}[!tb]
\centering
\caption{NAFNet, adapted protocol over one leave one subject out fold per subject, against the no
adaptation protocol of Table 3 of the main paper on the same checkpoint. Clinical rows are percentage
change from backbone output to corrected output. The spread is the standard deviation across
subjects, the count is the number of subjects on which the measure improved, and $p$ is a two sided
Wilcoxon signed rank test on the paired per subject values.}
\label{tab:supp_adapted}
\footnotesize
\begin{tabular}{lrrrr}
\toprule
Measure & No adaptation & Adapted & Improved & $p$ \\
\midrule
\multicolumn{5}{l}{\textit{Duke17, 16 subjects}} \\
$\Delta$PSNR (dB) & $-0.17$ & $-0.23 \pm 0.44$ & 6 of 16 & 0.093 \\
$\Delta$CNR & $+7.51$ & $+10.68 \pm 1.55$ & 16 of 16 & $3.1{\times}10^{-5}$ \\
$\Delta$TCI & $+4.00$ & $+2.69 \pm 4.16$ & 12 of 16 & 0.018 \\
$\Delta$EPI & $+1.37$ & $+1.49 \pm 0.78$ & 15 of 16 & $6.1{\times}10^{-5}$ \\
$\Delta$BS & $+4.15$ & $+4.13 \pm 5.18$ & 13 of 16 & 0.0017 \\
$\Delta$ENL & $+21.35$ & $+24.29 \pm 17.72$ & 15 of 16 & $6.1{\times}10^{-5}$ \\
$\Delta$SNR & $+5.43$ & $+2.95 \pm 8.97$ & 9 of 16 & 0.50 \\
\midrule
\multicolumn{5}{l}{\textit{Duke2013, 18 subjects}} \\
$\Delta$PSNR (dB) & $-0.35$ & $-0.45 \pm 0.47$ & 4 of 18 & 0.0016 \\
$\Delta$CNR & $+7.37$ & $+15.43 \pm 4.09$ & 18 of 18 & $7.6{\times}10^{-6}$ \\
$\Delta$TCI & $+1.41$ & $+5.15 \pm 3.75$ & 17 of 18 & $1.5{\times}10^{-5}$ \\
$\Delta$EPI & $+0.85$ & $+0.90 \pm 0.89$ & 15 of 18 & 0.00084 \\
$\Delta$BS & $+0.51$ & $+2.96 \pm 2.97$ & 14 of 18 & 0.0010 \\
$\Delta$ENL & $+28.80$ & $+9.23 \pm 28.61$ & 8 of 18 & 0.80 \\
$\Delta$SNR & $+9.29$ & $-5.29 \pm 13.97$ & 3 of 18 & 0.010 \\
\bottomrule
\end{tabular}
\end{table}

Two things hold on both scanners. Adaptation raises contrast, by 3.2 points on Duke17 and 8.1 on
Duke2013, and every subject improves. Adaptation also costs fidelity and signal to noise ratio
relative to the frozen model, a further 0.06 and 0.10 dB, with one and two subjects losing more than
1 dB, and a signal to noise ratio 2.5 and 14.6 points lower. The mean absolute correction rises from
0.0104 to 0.0121 on Duke17 and from 0.0118 to 0.0164 on Duke2013. The other measures do not agree
between the two sets. On Duke17 the tissue contrast index falls from $+4.0$ to $+2.7$ percent while
the speckle measure rises, and no measure falls below zero, although the signal to noise ratio is no
longer significant. On Duke2013 the tissue contrast index, edge preservation and boundary sharpness
all rise, the speckle measure changes on eight subjects one way and ten the other and is not
significant, and the signal to noise ratio falls by 5.3 percent, a decrease on fifteen of eighteen
subjects that is significant at the one percent level.

We draw one conclusion. Adaptation on fifteen to seventeen images of the target scanner reliably
improves one measure, contrast, which the frozen model already improves on both sets. It pays for
that in fidelity and signal to noise ratio on both sets, its effect on the remaining measures
changes direction between two scanners, and it does so from a checkpoint selected on the held out
subject. The frozen model gives six positive measures on both sets with no fitting on the target
data, and it is the protocol the main paper reports and the one we recommend.
